% ID: SEL-D0013-BC-E05
% ESTADO: cotejada. Segunda pasada visual completa y borrador compilado revisados.
% Fuente: Selectividad/matematicas_aplicadas/matematicas_aplicadas_P_A.pdf
% SHA-256: 0305750f6b6c3cc3edbba71a19cd7cce4276935dff5a0b7d964b02eaedeaa222
% Página: 2; bloque C; ejercicio original 5.

Una urna contiene 6 bolas rojas y 4 azules. Se extrae una bola al azar y se
reemplaza por seis bolas del otro color. A continuación, se vuelve a extraer
una segunda bola de la urna.
\begin{enumerate}[label=\alph*)]
  \item \textbf{(1.5 puntos)} Calcule la probabilidad de que la segunda bola
  extraída sea roja.
  \item \textbf{(1 punto)} Si sabemos que la segunda bola extraída es azul,
  \textquestiondown{}cuál es la probabilidad de que también lo haya sido la primera?
\end{enumerate}
